\topic{Field comparison}{Pressure across formulations at $Re=400$}
\lessonfigure{field_p_re400_n64}{Retained Re=400 zero-mean pressure and computed velocity streamlines at h=1/64: 64 FV pressure cells per side and 65 boundary-inclusive FD nodes. The common RGB pressure scale spans -2 to 2, with symmetric-log normalization linear below absolute pressure 0.01. FD pressure is reconstructed from the converged velocity field; Ghia provides no pressure reference.}{field-p-re400-n64}
The cavity has unit side length and a unit-speed upper lid; the other walls are
stationary. Pressure is nondimensionalized by $\rho U_{\rm lid}^2$ and reported
with zero spatial mean, as in Figure~\ref{fig:field-p-re100}.
The common color limits make pressure differences comparable across Reynolds
numbers. The color mapping is nonlinear, so use the numerical colorbar labels
rather than color distance to judge magnitude. Black curves follow velocity.

\textbf{Read the physics.} Locate the low-pressure vortex core and the pressure
rise near the upper-right turning region. Explain why similar streamline
patterns do not by themselves validate reconstructed pressure. Use the
pressure-gradient discrepancy in the retained measurement table as a separate
consistency check; it is not a comparison with an exact pressure field.
\newpage
\topic{Field comparison}{Pressure across formulations at $Re=1000$}
\lessonfigure{field_p_re1000_n64}{Retained Re=1000 zero-mean pressure and computed velocity streamlines at h=1/64: 64 FV pressure cells per side and 65 boundary-inclusive FD nodes. The common RGB pressure scale spans -2 to 2, with symmetric-log normalization linear below absolute pressure 0.01. FD pressure is reconstructed from the converged velocity field; Ghia provides no pressure reference.}{field-p-re1000-n64}
Use Figures~\ref{fig:field-p-re100}, \ref{fig:field-p-re400-n64} and
\ref{fig:field-p-re1000-n64} to compare pressure at the same $h/L=1/64$.
The geometry, lid speed, pressure gauge and color scale are common. The
discontinuous lid velocity produces demanding corner gradients; changing the
plot palette does not improve their spatial resolution.

\textbf{Separate appearance from accuracy.} Identify a pressure feature that
changes with Reynolds number and one region where formulation differences
remain. Compare against the Ghia velocity profiles and native residuals already
provided, then explain why neither test establishes pressure accuracy. The
32-cell pressure panels in the notebook provide a distinct coarse-grid check.
